%% ─────────────────────────────────────────────────────────────────────────────
\section{Dataset and Preprocessing}
%% ─────────────────────────────────────────────────────────────────────────────

\subsection{Source Corpus: Ship of Theseus}

Our corpus is the \emph{Ship of Theseus} paraphrased dataset of Tripto et al.~\cite{tripto2024ship}.
It is organised along three axes:

\begin{itemize}[noitemsep,topsep=2pt]
  \item \textbf{Seven datasets (domains):} XSUM (BBC news), Yelp (reviews), CMV (Reddit
        ChangeMyView arguments), ELI5 (Reddit explanatory answers), SciGen (scientific writing),
        TLDR (Reddit summaries) and WP (Reddit WritingPrompts stories).
  \item \textbf{Seven sources} that wrote the original text: Human, and six LLMs (OpenAI, LLaMA,
        PaLM, Tsinghua, BigScience, Eleuther-AI) prompted with the first 30 tokens of the matching
        human text. Human and LLM texts from the same prompt share a document \texttt{key}.
  \item \textbf{Seven paraphraser configurations} of four models, each applied three times in
        sequence (T0$\to$T1$\to$T2$\to$T3): ChatGPT (GPT-3.5-turbo), PaLM2, Pegasus in a
        \emph{full} (all sentences) and \emph{slight} (25\% of sentences) setting, and
        DIPPER~\cite{krishna2023paraphrasing} at default, \emph{low} and \emph{high} lexical diversity.
\end{itemize}

Each dataset ships as two files. The \textbf{paraphrased} file holds every original with its 21
paraphrases (7 configurations $\times$ T1--T3), keyed by \texttt{source}, \texttt{key} and
\texttt{version\_name} (\texttt{original}, \texttt{chatgpt}, \texttt{chatgpt\_chatgpt}, \ldots).
The \textbf{train} file holds a disjoint set of un-paraphrased texts from the same seven sources.
We use the paraphrased files for decay analysis (RQ1), paraphraser identification (RQ3) and as the
attribution test set (RQ2), and the train files to train the attribution model (RQ2).

As a controlled reference we keep our Phase-1 rule-based \textbf{WordNet} simulator, which
replaces 25\%/38\%/50\% of content words with same-POS WordNet synonyms at T1/T2/T3.

\subsection{Preprocessing Pipeline}

All 14 files are processed by one logged pipeline (\texttt{src/data/preprocess.py}):

\noindent\textbf{1. Cleaning.} HTML entities are decoded, Unicode is NFKC-normalised, curly quotes
and dashes are mapped to ASCII, and whitespace is collapsed. Sources and paraphrasers format text
differently (e.g.\ DIPPER outputs begin with a space; some models emit curly quotes), so without
this step punctuation and tokenisation features would measure formatting rather than style.
About 44\% of paraphrased rows were modified; no words or punctuation marks are removed.
We deliberately do \emph{not} lower-case, remove stopwords, stem, or truncate: these would
destroy the stylistic signal under study, and truncating T0 would misalign it with paraphrases
produced from the full document.

\noindent\textbf{2. Chain alignment.} Rows are pivoted to one row per (source, key) holding T0 and
all 21 paraphrases. Only documents with a complete chain for \emph{all seven} configurations are
kept, so every paraphraser is compared on identical texts. Retention is 90--97\% for most
domains but 58\% for SciGen and 69\% for WP, where many PaLM chains are missing.

\noindent\textbf{3. Length filter.} Documents whose T0 has fewer than 30 words are removed
($\approx$2\%), since ratio features and BLEU are unstable on very short texts.

\noindent\textbf{4. Sanity checks.} For each configuration we count unchanged paraphrase steps
(mostly Pegasus-slight and PaLM; kept, as they are genuine model behaviour) and T1 outputs sharing
fewer than 10\% of their words with T0. The latter flags \textbf{DIPPER (high)} in five of seven
domains (54--95\% of outputs): its released chains are largely unrelated text (e.g.\ narrative
fiction in place of a news article).

\noindent\textbf{5. Leakage-free splits.} The train files are split 85/15 into train/validation
\emph{by key}, so a human text and the LLM texts generated from its prompt never straddle the
split. We verified that the paraphrased and train files share no keys, so RQ2 test documents are
never seen in training.

\noindent\textbf{Output.} 16{,}386 aligned chains (2{,}269 Human) and 19{,}200 training texts.
For the analyses in this update we sample \textbf{100 Human documents per domain} (700 in total,
seed 42), since RQ1 concerns the decay of a human author's identity.
Table~\ref{tab:corpus_stats} summarises this sample.

\begin{table}[h]
\small
\centering
\caption{Human T0 originals used in this update (100 per domain, after cleaning and the
30-word filter). Vocabulary counts lower-cased alphabetic tokens.}
\label{tab:corpus_stats}
\begin{tabular}{lrrrr}
\toprule
\textbf{Domain} & \textbf{Docs} & \textbf{Avg.\ words} & \textbf{Avg.\ sents} & \textbf{Vocab} \\
\midrule
XSUM (news)       & 100 & 300.0 & 15.1 & 5{,}952 \\
Yelp (review)     & 100 & 131.5 &  9.4 & 2{,}677 \\
CMV (opinion)     & 100 & 401.8 & 23.9 & 4{,}899 \\
ELI5 (qa)         & 100 & 242.2 & 12.5 & 4{,}126 \\
SciGen (science)  & 100 & 163.0 &  7.0 & 3{,}181 \\
TLDR (summary)    & 100 &  84.7 &  4.8 & 2{,}458 \\
WP (story)        & 100 & 418.3 & 33.7 & 6{,}280 \\
\midrule
Total             & 700 & 248.8 & 15.2 & 15{,}888 \\
\bottomrule
\end{tabular}
\end{table}
